\begin{textsample}{POS Dim 3 – persona\_gemini – Score 17.00 – t777\_gemini.txt}  \label{ex:f3_pos_050}
With the launch of the S9 smartphone, Samsung is falling behind rivals like Apple, Google, and Facebook.
These \textbf{companies} have all committed to 100 % renewable \textbf{energy} and are making progress toward that goal.
Currently, Samsung \textbf{uses} only 1 % renewable \textbf{energy}.
This is a significant \textbf{problem} because 80 % of a gadget 's carbon \textbf{emissions} happen during manufacturing, which is largely powered by fossil fuels like coal and \textbf{gas}.
The \textbf{scale} of this \textbf{energy} \textbf{consumption} is vast ; the \textbf{electricity} \textbf{used} for smartphone \textbf{production} between 2007 and 2017 is \textbf{equivalent} to Japan 's annual \textbf{electricity} \textbf{consumption}.
Last year, 1.53 \textbf{billion} smartphones were sold, leading to a massive \textbf{waste} of resources and a growing e-waste crisis. \textbf{Customers} are ready for a change.
A 2015 \textbf{study} showed that 66 % of \textbf{customers} are willing to pay more for \textbf{products} from sustainable \textbf{brands}.
In a recent Greenpeace survey, 85 % of over 8,000 supporters said they would be more likely to buy Samsung \textbf{products} if the \textbf{company} was a leader in sustainability.
It 's \textbf{time} to tell Samsung to switch to renewables.
Post a comment on Samsung ’s Facebook page and demand they commit to 100 % renewable \textbf{energy}.

% matched lemmas: billion, brand, company, consumption, customer, electricity, emission, energy, equivalent, gas, problem, product, production, scale, study, time, use, waste
\end{textsample}
